\section{The obstruction from list distinctness}\label{sec:obstruction}
Fix a finite field \(k\). Consider the test that accepts a list precisely when
its records are pairwise distinct. Sorting gives one way to perform the
test. Our lower bound applies to every tensor network computing it, even
when a single network tests many lists at once. The contradiction will
come from estimating the cost per input bit. A matrix multiplication
exponent below \(9/4\) would drive that cost below the lower bound after
a fixed number of compressions.

\subsection{Two estimates for networks on bits}
A \emph{bit tensor} of arity \(u\) is a function \(\{0,1\}^u\to k\).
In a tensor network each vertex carries such a table. A wire pairs two
argument positions, and each unpaired position is a named exterior bit.
For a fixed setting of the exterior bits, multiply the vertex entries
and sum over all settings of the internal wires. This gives the tensor
represented by the network. A network is \emph{trivalent} if each
vertex has arity at most three. Its \emph{size} counts the vertices and
exterior bits. Its \emph{volume} counts the vertices and all wires,
including the exterior wires. These counts agree up to an absolute
factor by trivalency. For networks of unrestricted arity we count
vertices plus their total arity.

The first estimate starts with an arbitrary trivalent network. Its
conclusion gives a network for a tensor power, with separate exterior
bits for each copy. On these bits it must compute the product of the
original answers.

\begin{theorem}[Compression of tensor powers]\label{thm:shaving}
Suppose \(\omega_k<9/4\). For every \(\eta>0\) there is a \(\gamma>0\)
such that the following holds for all sufficiently large \(n\).
Write \(h=\lceil\log_2 n\rceil\). If a trivalent \(k\)-network of size
\(m\le n^{1.05}\) with \(e\ge n\) exterior bits represents \(V\), then
for some integer \(d\ge1\) its power \(V^{\otimes d}\) has a trivalent
\(k\)-network with \(de\) distinct exterior bits and size at most
\begin{equation}\label{eq:shave}
                   d\big(mh^{-\gamma}+eh^\eta\big).
\end{equation}
Here \(d\le n^{o(1)}\) uniformly in the input network for fixed \(k,\eta\).
\end{theorem}

Section~\ref{sec:shaving} proves Theorem~\ref{thm:shaving}. Per copy,
the first term saves a power of \(h\) on the original size. The second
charges \(h^\eta\) for each exterior bit. The choice of \(\eta\) lets
us compare this charge with the cost of list distinctness.

For \(w\ge1\), let \(D_w\) be the zero-one tensor for a list of \(2^w\)
words of \(2w\) bits each. Its value is one exactly when the words
are pairwise distinct. It has \(n_w=2w\,2^w\) exterior bits. The tensor
\(D_w^{\otimes t}\) tests distinctness separately within each of \(t\)
independently specified lists. The network may connect vertices across
lists. Let \(\operatorname{vol}_k\) denote minimum trivalent volume over \(k\).

\begin{theorem}[Volume of list distinctness]\label{thm:many}
There is a constant \(\sigma>0\), depending only on \(k\), such that for
all sufficiently large \(w\) and every integer \(1\le t\le 2^w\),
\begin{equation}\label{eq:many}
                  \operatorname{vol}_k(D_w^{\otimes t})
                          >t n_w w^\sigma .
\end{equation}
\end{theorem}

Section~\ref{sec:list} proves Theorem~\ref{thm:many}. Its uniformity
in \(t\) is essential because the replication factor in
Theorem~\ref{thm:shaving} is chosen by the compression.

\subsection{Comparison at a fixed input length}
\begin{proposition}\label{prop:finitefield}
Over every finite field \(k\), \(\omega_k\ge 9/4\).
\end{proposition}
\begin{proof}
A sorting network on \(2^w\) records uses \(O(2^w w^2)\) comparisons
\citep{batcher}. Comparing and conditionally exchanging two \(2w\)-bit
words takes \(O(w)\) Boolean gates. After sorting, the conjunction of
the inequalities between adjacent words computes \(D_w\). To realize
this circuit as a tensor network, use the zero-one table of each gate,
copy bits at fan-outs with equality tables, and pin the output to one.
For each accepted input there is exactly one consistent assignment of
the internal wires. This gives a trivalent network over \(k\) of size
\(m_0=O(n_w w^2)\).

Suppose \(\omega_k<9/4\). Choose \(0<\eta<\sigma\), take \(\gamma\)
from Theorem~\ref{thm:shaving}, and fix an integer \(J>2/\gamma\).
Set \(n=n_w\) and \(h=\lceil\log_2 n\rceil=w+O(\log w)\). These
parameters stay fixed through all \(J\) applications of the theorem.
At step \(j\) write \(t_j\) for the number of lists, \(m_j\) for the
network size and \(x_j=m_j/(t_jn)\) for the size per exterior bit.
Initially \(t_0=1\) and \(x_0=O(w^2)\). With replication factor \(d_j\),
\eqref{eq:shave} gives
\[
         t_{j+1}=d_jt_j,\qquad
         x_{j+1}\le h^{-\gamma}x_j+h^\eta .
\]
Inductively, at each of the first \(J\) steps the network has at least
\(n\) exterior bits and \(t_j=n^{o(1)}\), since the bound on \(d_j\)
is uniform and \(J\) is fixed. The recurrence gives
\(x_j=O(w^2+h^\eta)\), so \(m_j=t_j n x_j\le n^{1.05}\) for large \(w\).
This verifies the size hypothesis at each step. Summing the recurrence gives
\[
       x_J \le h^{-J\gamma}x_0+
                     h^\eta\sum_{r=0}^{J-1}h^{-r\gamma}
               = O(w^\eta).
\]
Since \(\log_2 t_J=o(w)\), Theorem~\ref{thm:many} applies to these
\(t_J\) lists. Their network has volume \(O(t_J n w^\eta)\) by
trivalency, which violates \eqref{eq:many} for sufficiently large \(w\).
\end{proof}

\subsection{Passing from finite fields to arbitrary fields}
The step that permits specialization is to extract a single bilinear
identity with a strict rank saving. Its coefficients generate a
finitely generated ring, and its equations survive reduction to a finite field.

\begin{lemma}[Bilinearization of a formal circuit]\label{lem:bilinearize}
If a formal circuit for multiplying two \(s\)-by-\(s\) matrices over
\(F\) uses \(q\) nonscalar multiplication or division gates, then
\begin{equation}\label{eq:gate-rank}
                     \operatorname{rank}_{\overline F}(T_s)\le 2q .
\end{equation}
Here \(q\) counts a multiplication when both inputs are nonconstant and
a division when its divisor is nonconstant.
\end{lemma}
\begin{proof}
Over the infinite field \(\overline F\), choose matrices \(A,B\) at
which all intermediate rational functions are regular and every divisor
is nonzero. Evaluate the circuit on \(A+uX,B+vY\) with \(u^2=v^2=0\).
Each gate has a unique expansion
\[
       f=f_{00}+u f_{10}+v f_{01}+uv f_{11}.
\]
Here \(f_{00}\) is a scalar, \(f_{10}\) is linear in \(X\), and
\(f_{01}\) is linear in \(Y\). At a product gate \(f=ab\), the mixed
coefficient satisfies
\[
        f_{11}=a_{00}b_{11}+a_{11}b_{00}
                         +a_{10}b_{01}+a_{01}b_{10}.
\]
Only the last two terms require new products of linear forms.
At a quotient gate \(f=a/b\), use \(bf=a\) and \(b_{00}\ne0\) to obtain
\[
        b_{00} f_{11}=a_{11}-b_{11}f_{00}-b_{10}f_{01}-b_{01}f_{10}.
\]
This again requires at most two new products. Linear operations,
including multiplication or division by a field scalar, contribute none.
The mixed coefficients of the outputs are the entries of \(XY\).
They all belong to the span of at most \(2q\) products of an
\(X\)-linear form and a \(Y\)-linear form, giving \eqref{eq:gate-rank}.
\end{proof}

\begin{lemma}[Finite specialization]\label{lem:specialize}
If \(\omega_F<9/4\) over a field \(F\), then \(\omega_k<9/4\)
over some finite field \(k\).
\end{lemma}
\begin{proof}
The circuit bound has the form \(O(s^{9/4-\epsilon})\) for some fixed
\(\epsilon>0\). By Lemma~\ref{lem:bilinearize}, at a sufficiently large
side \(s\ge2\) we have a rank decomposition of \(T_s\) over \(\overline F\)
with \(r<s^{9/4}\) terms. Let
\(R\subseteq\overline F\) be the unital ring generated over \(\mathbb Z\)
by all of its coefficients. The decomposition identities are polynomial
equations over \(R\) with target coefficients zero and one.

Choose a maximal ideal of \(R\). Its residue field \(k\) is finitely
generated as a \(\mathbb Z\)-algebra. In positive characteristic Zariski's
lemma makes \(k\) a finite extension of its prime field. In characteristic
zero the same lemma over \(\mathbb Q\) would make \(k\) a number field.
Multiplying its algebra generators by a common positive integer \(M\)
would then make them algebraic integers. Consequently the generated
ring would lie in the ring of integers of \(k\) with \(M\) inverted.
This ring omits \(1/\ell\) for every prime \(\ell\nmid M\), whereas
the field \(k\) contains all of \(\mathbb Q\). Thus \(k\) is finite.
Reducing the identities gives \(\operatorname{rank}_k(T_s)\le r\).

Finally, take the \(j\)-th power of this fixed bilinear algorithm.
Each of the three linear maps can be applied one tensor coordinate at
a time. Since \(s^2\le r\) by flattening rank, every intermediate
array has at most \(r^j\) entries. The linear work and the \(r^j\)
scalar products together cost \(O_{r,s}(j r^j)\) operations at side
\(s^j\). Padding to the next power of \(s\) gives
\(\omega_k\le\log_s r<9/4\).
\end{proof}

\begin{proof}[Proof of Theorem~\ref{thm:square}]
Proposition~\ref{prop:finitefield} and Lemma~\ref{lem:specialize} give
the lower bound over every field. The upper bound over every field is
due to Amanov, van den Berg and Zuiddam \citep{abvz}.
\end{proof}
